\documentclass[oneside,a4paper,14pt]{extarticle}
\usepackage[a4paper,top=20mm,bottom=20mm,left=30mm,right=10mm]{geometry}
\usepackage{amsmath, amssymb, amsfonts, amsthm}
\usepackage[T2A]{fontenc}
\usepackage[utf8]{inputenc}
\usepackage[russian]{babel}
\usepackage{textcomp}
\usepackage{indentfirst}
\usepackage{graphicx}
\usepackage{float}
\usepackage{wrapfig}
\usepackage{caption}
\usepackage{tocloft}
\usepackage{seqsplit}
\allowbreak
\renewcommand{\cftsecleader}{\cftdotfill{\cftdotsep}}
\usepackage[breaklinks]{hyperref}
\usepackage{titlesec}
\usepackage{fancyvrb}
\usepackage{multirow}
\usepackage{xcolor}

\titleformat{\section}{\normalsize\bfseries}{\thesection}{1em}{}
\titleformat{\subsection}{\normalsize\bfseries}{\thesubsection}{1em}{}
\renewcommand\baselinestretch{1.33}\normalsize
\setlength{\parindent}{1.25cm}

\captionsetup[figure]{labelsep=period, name=Рис.}

\begin{document}
\newpage\thispagestyle{empty}
\newpage
\setcounter{page}{1}
\tableofcontents
\newpage

\section*{Цель}
\addcontentsline{toc}{section}{Цель}
Разработать и исследовать вычислительное устройство.

\section*{Задание}
\addcontentsline{toc}{section}{Задание}
Выполнить сдвиг 8-разрядного двоичного кода $X$ на $N$ разрядов в сторону старших разрядов (сдвиг логический).

В табл. 1 приведено распределение ячеек ОЗУ. В ячейке с адресом 00 располагается код операции (COP), в ячейке 01 — исходное число $X$, в ячейке 02 — количество позиций сдвига $N$, а в ячейке 03 — результат выполнения операции $Y$. В ячейке 00 ПНА записан начальный адрес микропрограммы сдвига (`0A`).

\begin{flushright}
	Таблица 1
\end{flushright}
\begin{center}
	Распределение ячеек ОЗУ
\end{center}
\begin{table}[H]
	\centering
	\begin{tabular}{|c|c|c|l|}
		\hline
		Адрес & Код  & Мнемоника & \multicolumn{1}{c|}{Комментарий} \\ \hline
		00    & 0000 & COP       & Код операции сдвига              \\ \hline
		01    & 0022 & X         & Исходное число                   \\ \hline
		02    & 0002 & N         & Количество позиций сдвига        \\ \hline
		03    & 0088 & Y         & Результат сдвига                 \\ \hline
	\end{tabular}
\end{table}

На рис. 1 изображена граф-схема основной микропрограммы. В ней производится выборка кода операции из ОЗУ, установка адресов операндов с помощью регистра-указателя $RF$, считывание исходного числа $X$ (в регистр $R0$) и количества позиций сдвига $N$ (в регистр $R9$), а также передача управления микропрограмме выполнения логического сдвига через таблицу ПНА (на адрес `0A`). После завершения операции управление возвращается на МК 07, и полученный результат из $R0$ записывается в ячейку 03 ОЗУ (МК 08).

\begin{figure}[H]
	\centering
	\includegraphics[height=0.7\textheight, keepaspectratio]{Полный.png}
	\caption{Граф-схема основной микропрограммы}
	\label{fig:main_mp}
\end{figure}

На рис. 2 приведена граф-схема второй микропрограммы — микропрограммы логического сдвига, реализующей обработку исходного числа в регистрах операционного устройства.

В начале выполнения микропрограммы (МК 0A) проверяется значение количества позиций сдвига $N$ в счётчике $R9$. Если $N = 0$ (признак $Z = 1$), цикл сдвига пропускается и осуществляется условный переход на МК 07 для записи результата в ОЗУ.

Если $N \neq 0$ ($Z = 0$), в МК 0B выполняется логический сдвиг разрядов числа в регистре $R0$ влево в сторону старших разрядов ($R0 := 2R0$). После каждого сдвига в МК 0C значение счётчика $R9$ уменьшается на единицу ($R9 := R9 - 1$), после чего выполняется безусловный переход обратно на проверку условия в МК 0A. Проверка счётчика повторяется до тех пор, пока его значение не станет равным нулю.

\begin{figure}[H]
	\centering
	\includegraphics[height=0.55\textheight, keepaspectratio]{Микропрограмма.png}
	\caption{Граф-схема микропрограммы логического сдвига}
	\label{fig:shift_mp}
\end{figure}

\newpage

\section*{Кодирование микропрограммы}
\addcontentsline{toc}{section}{Кодирование микропрограммы}

В табл. 2 приведено кодирование микропрограммы логического сдвига. В ней указаны коды микрокоманд, значения управляющих сигналов, а также коды микрокоманд для управления условными переходами и циклическим выполнением сдвига.

\begin{flushright}
	Таблица 2
\end{flushright}
\begin{center}
	Кодирование микропрограммы логического сдвига
\end{center}
\begin{table}[H]
	\centering
	\renewcommand{\arraystretch}{1.3}
	\resizebox{\textwidth}{!}{
		\begin{tabular}{|c|c|c|c|c|c|c|c|c|c|c|c|c|c|c|c|c|c|}
			\hline
			№   & МИ    & \multicolumn{2}{c|}{РЗУ} & \multicolumn{3}{c|}{Упр. АЛУ} & \multicolumn{3}{c|}{Упр. ОЗУ} & Шина                & МИ & \multicolumn{3}{c|}{Упр. усл.} & \multicolumn{3}{c|}{Упр. УУ}                                                                                                                        \\ \hline
			МК  & I8-0  & A                        & B                             & C0                            & \textasciicircum OE & SC & \textasciicircum CS            & \textasciicircum W           & \textasciicircum EA & D11-0 & I3-0 & A  & U & \textasciicircum CCE & C0 & \textasciicircum RLD & \textasciicircum OE \\ \hline
			00: & 343   & 0                        & F                             & 0                             & 1                   & 00 & 1                              & 1                            & 1                   & 000   & E    & 00 & 0 & 0                    & 1  & 1                    & 0                   \\ \hline
			01: & 303   & 0                        & F                             & 1                             & 0                   & 00 & 1                              & 1                            & 0                   & 000   & E    & 00 & 0 & 0                    & 1  & 1                    & 0                   \\ \hline
			02: & 337   & 0                        & 0                             & 0                             & 1                   & 00 & 0                              & 1                            & 1                   & 000   & E    & 00 & 0 & 0                    & 1  & 1                    & 0                   \\ \hline
			03: & 303   & 0                        & F                             & 1                             & 0                   & 00 & 1                              & 1                            & 0                   & 000   & E    & 00 & 0 & 0                    & 1  & 1                    & 0                   \\ \hline
			04: & 337   & 0                        & 9                             & 0                             & 1                   & 00 & 0                              & 1                            & 1                   & 000   & E    & 00 & 0 & 0                    & 1  & 1                    & 0                   \\ \hline
			05: & 143   & 0                        & E                             & 0                             & 0                   & 00 & 1                              & 1                            & 0                   & 000   & E    & 00 & 0 & 0                    & 1  & 1                    & 0                   \\ \hline
			06: & 237   & 0                        & E                             & 0                             & 1                   & 00 & 0                              & 1                            & 1                   & 000   & 2    & 00 & 0 & 0                    & 1  & 1                    & 0                   \\ \hline
			07: & 303   & 0                        & F                             & 1                             & 0                   & 00 & 1                              & 1                            & 0                   & 000   & E    & 00 & 0 & 0                    & 1  & 1                    & 0                   \\ \hline
			08: & 133   & 0                        & 0                             & 0                             & 0                   & 00 & 0                              & 0                            & 1                   & 000   & 0    & 00 & 0 & 1                    & 1  & 1                    & 0                   \\ \hline
			09: & \dots & .                        & .                             & .                             & .                   & .. & .                              & .                            & .                   & \dots & .    & .. & . & .                    & .  & .                    & .                   \\ \hline
			0A: & 133   & 0                        & 9                             & 0                             & 1                   & 00 & 1                              & 1                            & 1                   & 007   & 3    & 00 & 1 & 0                    & 1  & 1                    & 0                   \\ \hline
			0B: & 733   & 0                        & 0                             & 0                             & 1                   & 00 & 1                              & 1                            & 1                   & 000   & E    & 00 & 0 & 0                    & 1  & 1                    & 0                   \\ \hline
			0C: & 313   & 0                        & 9                             & 0                             & 1                   & 00 & 1                              & 1                            & 1                   & 00A   & 3    & 00 & 0 & 1                    & 1  & 1                    & 0                   \\ \hline
		\end{tabular}
	}
	\label{tab:microcommands}
\end{table}

\newpage

\section*{Экранные формы}
\addcontentsline{toc}{section}{Экранные формы}

В ходе выполнения лабораторной работы в имитаторе микропрограммируемого микропроцессора были получены экранные формы ОУ и УУ (рис. \ref{fig:ou_uu}), а также экранные формы ОЗУ, ПНА и ПА (рис. \ref{fig:ozu_pna}).

\begin{figure}[H]
	\centering
	\includegraphics[width=0.95\textwidth]{окно.png}
	\caption{Экранная форма ОУ и УУ}
	\label{fig:ou_uu}
\end{figure}

\begin{figure}[H]
	\centering
	\includegraphics[width=0.95\textwidth]{память.png}
	\caption{Экранная форма ОЗУ, ПНА и ПА}
	\label{fig:ozu_pna}
\end{figure}

\newpage

\section*{Вывод}
\addcontentsline{toc}{section}{Вывод}
В ходе лабораторной работы было разработано вычислительное устройство для выполнения логического сдвига 8-разрядного двоичного кода $X$ на $N$ позиций в сторону старших разрядов. Было выполнено распределение регистров и ячеек ОЗУ, построены граф-схемы основной микропрограммы и микропрограммы логического сдвига.

В разработанной микропрограмме реализованы последовательное считывание операндов из ОЗУ с использованием регистра-указателя $RF$, дешифрация кода операции через таблицу начальных адресов ПНА, циклическое выполнение логического сдвига, контроль количества выполненных операций по счётчику $R9$ и запись полученного результата обратно в ОЗУ.

Работоспособность микропрограммы исследована в имитаторе микропрограммируемого микропроцессора.

\end{document}